
\documentclass[11pt,a4paper]{article}
\usepackage[utf8]{inputenc}
\usepackage[T2A]{fontenc}
\usepackage[russian,english]{babel}
\usepackage{amsmath,amssymb,amsthm,amsfonts}
\usepackage{geometry,booktabs,hyperref}
\geometry{margin=2.2cm}

\newtheorem{theorem}{Theorem}[section]
\newtheorem{lemma}[theorem]{Lemma}
\newtheorem{definition}[theorem]{Definition}
\newtheorem{corollary}[theorem]{Corollary}
\newtheorem{remark}[theorem]{Remark}
\newtheorem{axiom}{Axiom}[section]

\newcommand{\subfact}{\mathord{!}\,}
\newcommand{\Sgr}{\mathcal{S}}

\title{\textbf{Вероятностно-комбинаторная семизначная логика времени:\\
Тернарно-тетрарный маятник со стрелой (v7.1)}\\
\large Каноническая формализация: ортогональная дуальность $3 \leftrightarrow 4$,
каскад кодов Хэмминга и доказанная стрела времени}
\author{Алекс Мохов (математическое ядро верифицировано и расширено)}
\date{15.11.2026 (v7.1)}

\begin{document}
\maketitle

\begin{abstract}
Формулируется беспараметрическая дискретная теория времени. Из перестановочной
комбинаторики симметрических групп и теории кодирования выводится: (i)
арифметическая уникальность базиса $\{3,4\}$; (ii) семизначная алгебра истинности
на плоскости Фано $\mathrm{PG}(2,2)$; (iii) двухтактный маятник ролей
$3 \leftrightarrow 4$ с периодом $T_0 = 12 = |A_4|$; (iv) --- впервые в редакции
v7.1 --- \emph{доказанная} стрела времени: нарушение детального баланса каналов
удержания $\chi = W_1(4 \to 3)/W_1(3 \to 4) = 16/15 \neq 1$ выводится из точной
структуры провалов, а не постулируется. Все числовые утверждения подтверждены
полным перечислением и Монте-Карло-симуляцией ($2\cdot10^5$ циклов).
\end{abstract}

\part{Фундаментальная теория}

\section{Комбинаторный генезис}
Для $\Sgr_n$ числа встреч $D_{n,j} = \binom{n}{j}\subfact{n-j}$ подсчитывают
перестановки ровно с $j$ неподвижными точками; $\subfact{n} = D_{n,0}$,
$\Delta_n = n! - \subfact{n}$, $P_1(n) = \Delta_n/n! \to 1 - e^{-1}$.

\begin{axiom}[Тернарно-тетрарное расщепление]
Регистр состояния $R_7 = L_4 \oplus L_3$: тетрарная логика $L_4$ — данные
(4 информационных бита), тернарная логика $L_3$ — контроль (3 проверочных бита,
основание $q = 3$). Время — длительность удержания текущего распределения ролей
до акта инверсии $L_4 \leftrightarrow L_3$.
\end{axiom}

\begin{theorem}[Арифметическая уникальность $n = 3$]
$n = 3$ — единственный корень уравнения баланса
$\Delta_{n+1} - \Delta_n = n(n+1) - 1$ на $n \ge 3$; квант перехода
$\Delta_4 - \Delta_3 = 11 = \subfact{4} + \subfact{3}$.
\end{theorem}

\section{Маятник со стрелой (новое в v7.1)}
\begin{definition}[Вставка и стягивание]
Прямой канал: для $\pi \in M_1(3)$ и $i \in \{1,\dots,4\}$ полагаем
$\sigma = (i\ 4) \circ \pi$ (16 элементарных событий).
Обратный канал: для $\sigma \in M_1(4)$ полагаем $\tau = \sigma|_{\{1,2,3\}}$
(15 элементарных событий). Каналы удержания:
$W_1(3 \to 4) = 13/16$, $W_1(4 \to 3) = 13/15$.
\end{definition}

\begin{theorem}[Маяпник со стрелой]
$\chi = 16/15 \neq 1$; стационарные доли элементарных событий $16/31 : 15/31$;
квант асимметрии $3/16 - 2/15 = 13/240 > 0$.
\end{theorem}

\begin{lemma}[Структура провалов прямого канала]
$\sigma = (i\ 4)\circ\pi$ — полный беспорядок $\iff$ $\pi$ — транспозиция $(a\ b)$
и $(i\ 4)$ её не пересекает. Три провала — в точности три паросочетания
четвёрки: $\{12|34\}, \{13|24\}, \{14|23\}$.
\end{lemma}

\begin{lemma}[Структура провалов обратного канала]
$\tau$ — полный беспорядок на $\{1,2,3\}$ $\iff$ $\sigma|_{\{1,2,3\}}$ — 3-цикл
и $\sigma(4) = 4$ (ровно 2 случая).
\end{lemma}

\begin{proof}[Доказательство Теоремы 2]
Из лемм: потери $3/16$ и $2/15$; нарушение детального баланса
$\chi = 16/15$; доли элементарных событий пропорциональны числам событий каналов.
\end{proof}

\part{Верификация}
Полное перечисление: $W_1(3\to4)=13/16$, $W_1(4\to3)=13/15$, $\chi=16/15$,
$3/16-2/15=13/240$. Монте-Карло ($2\cdot10^5$ циклов $T_0=12$): эмпирика
$0.812623$ и $0.866737$; $E[\text{потерь/цикл}] = 1.925$;
$P(\text{сброс за цикл}) = 0.596747$ (пуассон-биномиаль, MC $0.595290$);
$\Delta\tau_{\text{metric}} = 0.441295$ нат; фазовый сдвиг $\ln(13/11) = 0.167054$.

\part{Следствия}
$\gamma = 3/12 = 1/4$; лоренцева сигнатура $(-,+,+,+)$; блокировка $n \ge 5$
($\subfact{5} = 44$ не является степенью простого).

\appendix
\section{Код верификации}
Полный исполняемый пакет Python 3 прилагается отдельно.

\end{document}
